\documentclass[11pt,a4paper]{article}
\usepackage{amsmath,amssymb,amsthm,mathtools}
\usepackage{geometry}
\usepackage{enumitem}
\usepackage{hyperref}
\geometry{top=25mm,bottom=25mm,left=25mm,right=25mm}

\newtheorem{definition}{定義}[section]
\newtheorem{theorem}[definition]{定理}
\newtheorem{proposition}[definition]{命題}
\newtheorem{lemma}[definition]{補題}
\newtheorem{remark}[definition]{注記}

\newcommand{\Ax}{\mathrm{Ax}}
\newcommand{\Pow}{\mathcal P}
\newcommand{\Solver}{\mathcal S}
\newcommand{\Zset}{\mathcal Z}
\newcommand{\Th}{\mathrm{Th}}
\newcommand{\Rstar}{\mathcal R^*}
\newcommand{\N}{\mathbb N}
\newcommand{\R}{\mathbb R}

\title{ECA における濃度構造の統合的到達点\\
\large ECA\_1--24 および既存の濃度関連文書の統合}
\author{}
\date{}

\begin{document}
\maketitle
\tableofcontents

\section{目的と位置づけ}

本稿では、これまでに構築された
\[
\mathrm{ECA}_1,\ldots,\mathrm{ECA}_{24}
\]
および既存の濃度関連文書における議論を、濃度追跡という観点から
一つの連続した構造として整理する。

本稿は新しい公理や追加条件を導入するものではない。既存文書で構築された
\[
\operatorname{Ax}(U)\longrightarrow
\Rstar\longrightarrow
\Pow(\Rstar)\longrightarrow
I\longrightarrow
\Solver_I\longrightarrow
\Zset
\]
という経路と、通常数学における
\[
\N\longrightarrow\Pow(\N)\longrightarrow\R
\]
という濃度構造を統合して記述する。

特に、公理候補集合、公理選択空間、Solver 集合、ZFC に対応する
Solver 集合 $\Zset$、Solver が表現する集合論的対象、および
ECA の評価値を同一視しない。

\section{ECA における基本構造}

ECA_6 以降で整理された構造では、ECA が扱う公理および公理スキームの全体を
\[
\Ax
\]
とし、公理選択空間を
\[
I\subseteq\Pow(\Ax)
\]
とする。

Solver は
\[
s=(\mathcal A_s,\mathcal R_s)
\]
であり、$\mathcal A_s$ は公理選択、$\mathcal R_s$ は推論規則および構成制約である。

\[
\Solver_I=
\left\{
s=(\mathcal A_s,\mathcal R_s)\mid
\mathcal A_s\in I
\right\}.
\]

公理選択復元写像は
\[
\mathfrak A:\Solver_I\to I,
\qquad
\mathfrak A(s)=\mathcal A_s
\]
である。

Solver が表現する理論を
\[
\mathcal T:\Solver_I\to\Th
\]
とし、
\[
\boxed{
\Zset=
\left\{
s\in\Solver_I\mid
\mathcal T(s)\cong_{\mathrm{Th}}\mathrm{ZFC}
\right\}
}
\]
と定める。

理論同型類を $[\mathrm{ZFC}]$ と書けば、
\[
\boxed{
\Zset=\mathcal T^{-1}([\mathrm{ZFC}])
}
\]
となる。

\section{数学的宇宙と公理候補集合}

ECA_12--15 では、数学的宇宙 $U$ とその公理候補集合
$\operatorname{Ax}(U)$ の関係を整理している。

\[
\operatorname{Ax}_{\mathrm{map}}:
\mathfrak U_{\mathrm{ECA}}\to\Pow(\Ax),
\qquad
\operatorname{Ax}_{\mathrm{map}}(U)=\operatorname{Ax}(U)
\]
を考え、
\[
\Ax=
\bigcup_{U\in\mathfrak U_{\mathrm{ECA}}}
\operatorname{Ax}(U)
\]
によって ECA 全体の公理候補集合へ接続する。

ECA_3 以来、ECA が扱う候補には個々の公理と公理スキームが含まれる。
ECA_2 では分離公理・置換公理が公理スキームとして扱われることも明示されている。

したがって、スキーム $\Sigma$ の個別インスタンス
$\Sigma(\varphi)$ を個々の公理候補として扱う構成が可能であり、
ECA_15、ECA_19 ではこの構成が明示されている。

\section{公理スキームから可算無限公理候補族へ}

可算個の入力
\[
\{\varphi_n\mid n\in\N\}
\]
に対して
\[
r_n=\Sigma(\varphi_n)
\]
とする。

各 $r_n$ が
\[
r_n\in\operatorname{Ax}(U)
\]
を満たし、かつ
\[
n\neq m\Longrightarrow r_n\neq r_m
\]
ならば、
\[
\Rstar=\{r_n\mid n\in\N\}
\subseteq\operatorname{Ax}(U)
\]
であり、
\[
\boxed{|\Rstar|=\aleph_0}
\]
を得る。

ここで重要なのは、$\Rstar$ 自体の基数と、その冪集合の基数を区別することである。

\section{可算無限族から冪集合へ}

\[
|\Rstar|=\aleph_0
\]
ならば、カントールの冪集合定理から
\[
|\Pow(\Rstar)|
=
2^{|\Rstar|}
=
2^{\aleph_0}.
\]

したがって
\[
\boxed{
\Rstar\longrightarrow\Pow(\Rstar)
}
\]
によって、可算無限の公理候補族から連続体濃度が現れる。

これは
\[
|\Rstar|=2^{\aleph_0}
\]
という主張ではない。連続体濃度を与えるのは
$\Pow(\Rstar)$ の側である。

\section{通常数学における連続体濃度}

通常数学では
\[
|\N|=\aleph_0
\]
であり、
\[
|\Pow(\N)|=2^{\aleph_0}.
\]

各 $A\subseteq\N$ に対して特性関数
\[
\chi_A:\N\to\{0,1\}
\]
を考え、
\[
x_A=
\sum_{n=1}^{\infty}\frac{\chi_A(n)}{2^n}
\]
とすれば、二進展開の非一意性を適切に処理することで
$\Pow(\N)$ と $[0,1]$ の間の全単射を構成できる。

さらに
\[
|[0,1]|=|\R|
\]
なので、
\[
\boxed{
|\R|=|\Pow(\N)|=2^{\aleph_0}
}
\]
となる。

またカントールの定理により
\[
\boxed{\aleph_0<2^{\aleph_0}}.
\]

\section{$\aleph_1$ と連続体濃度}

$\aleph_1$ を最初の非可算基数とする。
任意の基数 $\kappa$ について
\[
\aleph_0<\kappa\Longrightarrow\aleph_1\leq\kappa
\]
なので、
\[
\boxed{\aleph_1\leq2^{\aleph_0}}.
\]

従って通常数学では
\[
\boxed{
\aleph_0<\aleph_1\leq2^{\aleph_0}
}.
\]

CH は
\[
\boxed{2^{\aleph_0}=\aleph_1}
\]
であるから、既に成立している不等式に加えて
\[
2^{\aleph_0}\leq\aleph_1
\]
が必要となる。

\section{$\Rstar$ の部分集合族から公理選択空間へ}

ZFC を表現する基礎公理選択を
\[
\mathcal A_{\mathrm{ZFC}}
\]
とする。

任意の
\[
A\subseteq\Rstar
\]
に対して
\[
i_A=\mathcal A_{\mathrm{ZFC}}\cup A
\]
と定める。

既存文書で定式化された局所選択閉包条件
\[
\boxed{
\forall A\subseteq\Rstar,\qquad
\mathcal A_{\mathrm{ZFC}}\cup A\in I
}
\]
により、
\[
\Gamma:\Pow(\Rstar)\to I,
\qquad
\Gamma(A)=\mathcal A_{\mathrm{ZFC}}\cup A
\]
を得る。

さらに
\[
\Rstar\cap\mathcal A_{\mathrm{ZFC}}=\varnothing
\]
ならば $\Gamma$ は単射であるため、
\[
\Pow(\Rstar)\hookrightarrow I
\]
となり、
\[
\boxed{|I|\geq2^{\aleph_0}}.
\]

\section{公理選択から ZFC 保存 Solver へ}

ECA_23--24 では、固定した推論・構成規則 $\mathcal R_0$ の下で
\[
s_A=(i_A,\mathcal R_0)
\]
を構成する。

Stage 3 の条件は
\[
\boxed{
\forall A\subseteq\Rstar,\qquad
\mathcal T(s_A)\cong_{\mathrm{Th}}\mathrm{ZFC}
}.
\]

これは、各 $r_n$ が単独で ZFC と両立することではなく、
$\Rstar$ の任意の部分集合を追加しても ZFC と理論同型な理論を
表現できることを要求する。

基礎理論を
\[
T_0=
\mathcal T(\mathcal A_{\mathrm{ZFC}},\mathcal R_0)
\]
とし、
\[
\forall r\in\Rstar,\qquad T_0\vdash r
\]
という冗長性条件を候補とする。

同一言語・固定推論規則の設定でこれが成立すれば、任意の
$A\subseteq\Rstar$ に対して
\[
\operatorname{Th}(T_0\cup A)
=
\operatorname{Th}(T_0)
\]
という形の保存が得られる。

さらに
\[
T_0\cong_{\mathrm{Th}}\mathrm{ZFC}
\]
ならば、
\[
\forall A\subseteq\Rstar,\qquad s_A\in\Zset
\]
へ接続する。

\section{$\Pow(\Rstar)$ から $\Zset$ への単射}

以上を合わせ、
\[
|\Rstar|=\aleph_0,
\qquad
\Rstar\cap\mathcal A_{\mathrm{ZFC}}=\varnothing
\]
とし、選択可能性と ZFC 保存性が成立するとする。

\[
\Phi:\Pow(\Rstar)\to\Zset,
\qquad
\Phi(A)=s_A
\]
を定める。

$A\neq B$ なら
\[
\mathcal A_{\mathrm{ZFC}}\cup A
\neq
\mathcal A_{\mathrm{ZFC}}\cup B
\]
なので
\[
s_A\neq s_B.
\]
したがって $\Phi$ は単射であり、
\[
\boxed{
|\Zset|
\geq
|\Pow(\Rstar)|
=
2^{\aleph_0}
}.
\]

これは ECA_9--24 において段階的に構成されてきた
連続体濃度の $\Zset$ 内部への接続である。

\section{濃度追跡全体の統合}

以上の構造は
\[
\boxed{
\operatorname{Ax}(U)
\longrightarrow
\Rstar
\longrightarrow
\Pow(\Rstar)
\longrightarrow
I
\longrightarrow
\Solver_I
\longrightarrow
\Zset
}
\]
とまとめられる。

濃度については、
\[
|\Rstar|=\aleph_0,
\]
\[
|\Pow(\Rstar)|=2^{\aleph_0},
\]
\[
|I|\geq2^{\aleph_0},
\]
および ZFC 保存 Solver の構成条件が成立する場合に
\[
|\Zset|\geq2^{\aleph_0}
\]
となる。

さらに
\[
\Zset\subseteq\Solver_I
\]
なので
\[
|\Zset|\leq|\Solver_I|.
\]

従って、
\[
\boxed{
2^{\aleph_0}\leq|\Zset|\leq|\Solver_I|
}
\]
という濃度範囲が得られる。

ここで $|\Solver_I|\leq2^{\aleph_0}$ の一般上界は本構成からは確定していない。
従って $|\Zset|=2^{\aleph_0}$ を一般結論として同一視せず、
既存構成から得られる下界として扱う。

\section{ECA の評価値と通常の基数の分離}

ECA には
\[
S\longrightarrow i_S
\longrightarrow(D_S,C_S)
\longrightarrow
(E_{\mathrm d}(S),E_{\mathrm c}(S))
\longrightarrow
\mathcal G(S)
\]
という評価経路がある。

しかし、
\[
E_{\mathrm c}(S)\in[0,1]
\]
であることや、
\[
\Omega_{\mathrm{CH}}\subseteq[0,1]^4
\]
のような状態空間を考えられることから、通常の集合論における基数を
直接決定することはできない。

したがって、
\[
\boxed{
\text{ECA の型・評価値}
\neq
\text{通常の集合論における基数}
}
\]
を維持する。

同様に、
\[
|\R|=2^{\aleph_0}
\]
と
\[
|\Zset|\geq2^{\aleph_0}
\]
は異なる命題である。

前者は通常数学における実数集合の濃度、
後者は ECA の Solver 集合の一部として定義された $\Zset$ の
濃度の下界である。

\section{通常数学と ECA の二つの経路}

連続体濃度を扱う際には、二つの経路を区別する必要がある。

通常数学側では
\[
\N\longrightarrow\Pow(\N)\longrightarrow\R
\]
によって
\[
|\R|=2^{\aleph_0}
\]
を得る。

ECA 側では
\[
\operatorname{Ax}(U)
\longrightarrow
\Rstar
\longrightarrow
\Pow(\Rstar)
\longrightarrow
I
\longrightarrow
\Solver_I
\longrightarrow
\Zset
\]
によって、連続体濃度を公理選択および Solver 構造へ接続する。

この二つを混同せず、
\[
\boxed{
\text{通常数学上の濃度}
\quad\text{と}\quad
\text{ECA 上の Solver 構造}
}
\]
を別個に保持する。

\section{CH に対する現在の位置}

通常数学から
\[
\aleph_1\leq2^{\aleph_0}
\]
が得られている。

一方、CH のためには
\[
2^{\aleph_0}\leq\aleph_1
\]
が必要である。

既存の ECA 評価経路
\[
i_S
\longrightarrow
(D_S,C_S)
\longrightarrow
(E_{\mathrm d},E_{\mathrm c})
\longrightarrow
\mathcal G(S)
\]
だけから、この上界を導くことはできない。

従って CH を ECA から導出する場合に必要なのは、
\[
\boxed{
\Pow(\N)\hookrightarrow\omega_1
}
\]
に相当する構造である。

評価写像
\[
\Pow(\N)\to[0,1]
\]
や
\[
\Omega_{\mathrm{CH}}\subseteq[0,1]^4
\]
だけでは、この上界は得られない。

したがって、本稿で確認される連続体濃度の構造化と、
\[
2^{\aleph_0}=\aleph_1
\]
という CH の導出とは分離される。

\section{$\Rstar$ の位置づけ}

$\Rstar$ は
\[
|\Rstar|=\aleph_0
\]
を満たす可算無限の公理候補族である。

しかし、ECA における $\Rstar$ の役割はその基数値だけでは表せない。

\[
\Rstar
\longrightarrow
\Pow(\Rstar)
\longrightarrow
I
\longrightarrow
\Solver_I
\longrightarrow
\Zset
\]
という展開を通して、$\Rstar$ は連続体濃度を公理選択および Solver
構造へ接続する基礎となる。

したがって、
\[
\boxed{
|\Rstar|=\aleph_0,
\qquad
|\Pow(\Rstar)|=2^{\aleph_0}
}
\]
という二段階を維持する。

ここでいう「連続体濃度よりも複雑な濃度」という表現は、
$\Rstar$ 自身の基数が $2^{\aleph_0}$ より大きいという意味ではなく、
$\Rstar$ を起点として形成される多段の濃度・構造関係を、
単一の基数値だけでは表現できないという意味で理解する。

\section{統合された到達点}

以上を総合すると、これまでの ECA_1--24 と既存の濃度関連文書によって、
次の構造が構築されている。

\[
\boxed{
\operatorname{Ax}(U)
\longrightarrow
\Rstar
\longrightarrow
\Pow(\Rstar)
\longrightarrow
I
\longrightarrow
\Solver_I
\longrightarrow
\Zset
}
\]

その濃度関係は、
\[
\boxed{|\Rstar|=\aleph_0}
\]
から
\[
\boxed{|\Pow(\Rstar)|=2^{\aleph_0}}
\]
へ進み、選択可能性および ZFC 保存 Solver の構成を通じて
\[
\boxed{2^{\aleph_0}\leq|\Zset|}
\]
へ接続される。

また、
\[
\boxed{|\Zset|\leq|\Solver_I|}
\]
であるため、
\[
\boxed{
2^{\aleph_0}
\leq
|\Zset|
\leq
|\Solver_I|
}
\]
が現在の $\Zset$ に対する濃度範囲となる。

通常数学側では、
\[
\boxed{
|\R|=|\Pow(\N)|=2^{\aleph_0}
}
\]
および
\[
\boxed{
\aleph_0<\aleph_1\leq2^{\aleph_0}
}
\]
が確定している。

これに対して CH の最終条件は
\[
2^{\aleph_0}\leq\aleph_1
\]
であり、これは既存の ECA 評価値や連続型分類から自動的には得られない。

従って、現在までの ECA の濃度構造における到達点は、
CH の先取りではなく、
\[
\boxed{
\text{可算無限の公理候補族}
\longrightarrow
\text{その冪集合}
\longrightarrow
\text{連続体濃度}
\longrightarrow
\text{公理選択}
\longrightarrow
\text{Solver}
\longrightarrow
\mathcal Z
}
\]
という具体的な構造経路を形成し、連続体濃度を ECA の公理選択・Solver
構造へ接続するところまでを体系化したことにある。

\section{結語}

$\Rstar$ は可算無限の基礎族として構成され、その冪集合によって
\[
2^{\aleph_0}
\]
が現れる。

さらに、その部分集合族を公理選択へ接続し、ZFC 保存 Solver の族として
$\Zset$ に接続することで、
\[
2^{\aleph_0}\leq|\Zset|
\]
という連続体濃度の下界を構成できる。

この構造は、
\[
\operatorname{Ax}(U)
\to
\Rstar
\to
\Pow(\Rstar)
\to
I
\to
\Solver_I
\to
\Zset
\]
という一連の濃度追跡としてまとめられる。

同時に、
\[
\text{濃度}\neq\text{構造}
\]
および
\[
\text{ECA の評価値}\neq\text{通常の集合論における基数}
\]
という区別を維持する。

したがって、$\Rstar$ の意義は単独の可算濃度にあるのではなく、
その冪集合および後段の構造を通じて連続体濃度を ECA の内部構造へ
接続する点にある。

\end{document}
